% Part: first-order-logic
% Chapter: tableaux
% Section: rules-and-proofs

\documentclass[../../../include/open-logic-section]{subfiles}

\begin{document}

\iftag{FOL}
      {\olfileid[zh]{fol}{tab}{rul}}
      {\olfileid[zh]{pl}{tab}{rul}}

\olsection{\usetoken{P}{tableau}规则}

!!^a{tableau}是对!!a{sentence}在!!a{structure}中可以为真或为假的种种可能方式所做的一种系统性考察。表列的构件是!!{signed formula}：即!!{sentence}再加上一个真值「符号」，要么是 $\True$，要么是~$\False$。这些带符号的!!{formula}被排成一棵（向下生长的）树。

\begin{defn}
  一个 \emph{!!{signed formula}} 是这样一个二元组，它由一个真值与!!a{sentence}构成，即要么是：
  \[
  \sFmla{\True}{!A} \text{ 或 } \sFmla{\False}{!A}.
  \]
\end{defn}

直觉上，我们可以把 $\sFmla{\True}{!A}$ 读作「$!A$ 可能为
真」，把 $\sFmla{\False}{!A}$ 读作「$!A$ 可能为假」（在某个!!{structure}中）。

树中的每个!!{signed formula}要么是一个 \emph{假设}（它们列在树的最顶端），要么是由其上方的一个!!a{signed formula}通过若干推理规则之一得到的。
前一个!!{formula}的每个可能的!!{main operator}都有两条规则，一条用于符号为~$\True$ 的情形，
另一条用于符号为~$\False$ 的情形。有些规则允许树
分叉，有些规则则只向该支添加!!{signed formula}。
一条规则可以（而且往往必须）不是应用于紧接着其上面的!!{signed formula}，而是应用于
从根到该规则被应用的这一处之间的支上的任意!!{signed formula}。

当一支既包含 $\sFmla{\True}{!A}$
又包含 $\sFmla{\False}{!A}$ 时，该支是 \emph{闭的}。一个闭的!!{tableau}是这样一棵表列，其中每一支
都是闭的。按照直觉上的解释，任何一支都描述了一种
共同的（joint）可能性，但 $\sFmla{\True}{!A}$ 和 $\sFmla{\False}{!A}$
不可能共同地（jointly）成立。换言之，如果一支是闭的，那么它
所描述的那种可能性就被排除了。特别地，这意味着
一个闭的!!{tableau}排除了一切同时
使得每个形如 $\sFmla{\True}{!A}$ 的假设为真、每个
形如~$\sFmla{\False}{!A}$ 的假设为假的可能性。

一个\emph{针对 $!A$ 的}闭!!{tableau}是以
$\sFmla{\False}{!A}$ 为根~的一个闭!!{tableau}。如果存在这样的闭!!{tableau}，那么
~$!A$ 为假的一切可能性就都被排除了；即 $!A$
在每一个!!{structure}中都必然为真。

\end{document}
